\subsection{Trace Intervention Experiment}
\label{sec:intervention}

A central question in understanding LLM-based agent behavior is whether performance is driven primarily by the model's latent capabilities or by the conversational context accumulated during execution.
To disentangle these factors, we design a \emph{trace intervention} experiment that causally manipulates the agent's conversational history before it begins solving a task.

\subsubsection{Experimental Design}

\paragraph{Overview.}
The experiment proceeds in three phases: (1)~baseline data collection, (2)~trace registry construction, and (3)~intervention evaluation.
In the baseline phase, each agent solves every selected task multiple times without any intervention, producing a set of successful and failed conversation traces.
These traces are then organized into a registry that maps each task to pools of success and failure trajectories.
In the evaluation phase, partial conversation fragments from these traces are injected into the agent's message history \emph{before} it begins a new trial, and we measure the resulting change in task success rate.

\paragraph{Task selection.}
For each environment, we select tasks that exhibit \emph{mixed} success rates (strictly between 0\% and 100\%) in at least one of the two agent types tested (ReAct and Tool Calling).
A task qualifies for intervention if it has mixed results in \emph{either} agent type; once qualified, it is included for \emph{both} agent types to ensure equal task counts across conditions and enable direct comparisons.
This selection criterion targets the regime where intervention effects are most detectable---tasks that are neither trivially easy nor impossibly hard.

\paragraph{Trace collection and registry.}
The baseline phase runs 15 independent trials per task per agent, each with temperature $T{=}0.7$, producing stochastic variation across runs.
Every trial generates a complete conversation trace: a sequence of messages recording the agent's reasoning steps, tool invocations, and environment observations.
For each task, we partition these traces into two pools based on the trial outcome:
\begin{itemize}
    \item \textbf{Success pool}: traces from trials where the agent achieved the task objective.
    \item \textbf{Failure pool}: traces from trials where the agent did not achieve the task objective.
\end{itemize}
All traces within each pool are stored and shuffled; at each intervention trial, a trace is sampled uniformly at random from the appropriate pool.
This ensures that results are not confounded by properties of any single trace and that the agent encounters diverse reasoning trajectories across trials.

\paragraph{Intervention mechanism.}
The intervention operates via a \texttt{BEFORE\_TASK} hook in the agent framework.
Before the agent begins reasoning about a task, the hook:
\begin{enumerate}
    \item Selects a random trace from the appropriate pool (success or failure) for the current task, filtering to traces with a sufficient number of assistant steps for the requested intervention dosage.
    \item Extracts a prefix of the assistant's reasoning steps from that trace.
    \item Injects these steps into the agent's conversation history, along with the corresponding tool execution results.
\end{enumerate}

Crucially, the injected steps are not merely passive text: for \textbf{ReAct} agents, each reasoning step includes structured \texttt{<thought>}/\texttt{<action>}/\texttt{<action\_input>} segments, and the associated tools are re-executed against the live environment server, producing fresh observations that reflect the current task state.
For \textbf{Tool Calling} agents, the structured \texttt{tool\_calls} from the trace are replayed, with each function call executed and its result appended as a tool response message.
This faithful replay ensures that the environment state is consistent with the injected trajectory---the agent resumes from a state that actually corresponds to the injected reasoning, not a hallucinated one.

\paragraph{Intervention dosages.}
We parameterize the intervention by $n_s$, the number of assistant reasoning steps injected from the trace.
Two complementary indexing schemes are used:
\begin{itemize}
    \item \textbf{Forward indexing} ($n_s = 1, 2$): inject the first $n_s$ steps from the trace, providing the agent with the initial approach and early tool calls.
    \item \textbf{Backward indexing} ($n_s = n{-}2, n{-}1$): inject all steps except the last $|n_s|$, providing the agent with nearly the complete trajectory and leaving only the final steps to be completed autonomously.
\end{itemize}
This design creates a dosage gradient: at $n_s{=}1$, the agent receives minimal guidance (just the opening move); at $n_s{=}n{-}1$, it receives almost the entire trajectory and must only execute the final reasoning step and solution submission.
At each dosage level, we cross with two intervention types (success trace vs.\ failure trace), yielding $4 \times 2 = 8$ intervention conditions plus the unmodified baseline.

\paragraph{Experimental controls.}
Each condition is run with 15 trials per task using Claude Sonnet 4.5 via Amazon Bedrock, with temperature $T{=}0.7$ and workflow-level tool verbosity.
Separate environment server instances are maintained for each agent type to prevent state contamination between concurrent runs.
A run is skipped if a report from a prior execution already exists, preventing duplicate data collection.
All trials within a condition use identical agent configurations; only the injected trace content (or its absence) varies.

\subsubsection{Results: Resistor Network Environment}

We present the first complete set of results on the \textsc{Resistor Network} environment (level~1), which requires agents to reconstruct circuit topologies from pairwise resistance measurements between nodes.
This environment demands multi-step mathematical reasoning, systematic exploration of circuit configurations, and validation of proposed solutions against measured values.
Four tasks with mixed baseline success rates were selected; three had valid intervention traces for both success and failure pools.

\begin{table}[t]
\centering
\caption{Task success rates (\%) on the Resistor Network environment under trace intervention. \textbf{Baseline}: no injected trace. \textbf{Success/Failure}: partial traces from prior successful/failed runs injected before the agent begins. $n_s$: number of injected assistant steps; $n{-}k$ denotes all steps except the last $k$.}
\label{tab:intervention-resistor}
\small
\begin{tabular}{lcccccc}
\toprule
 & & \multicolumn{4}{c}{\textbf{Success Intervention}} \\
\cmidrule(lr){3-6}
\textbf{Agent} & \textbf{Baseline} & $n_s{=}1$ & $n_s{=}2$ & $n_s{=}n{-}2$ & $n_s{=}n{-}1$ \\
\midrule
ReAct         & 73.3 & 86.7 & 95.6 & 97.8 & \textbf{100.0} \\
Tool Calling  & 60.0 & 68.9 & 82.2 & 97.8 & \textbf{100.0} \\
\midrule
 & & \multicolumn{4}{c}{\textbf{Failure Intervention}} \\
\cmidrule(lr){3-6}
 & \textbf{Baseline} & $n_s{=}1$ & $n_s{=}2$ & $n_s{=}n{-}2$ & $n_s{=}n{-}1$ \\
\midrule
ReAct         & 73.3 & 48.9 & 55.6 & 11.1 & 11.1 \\
Tool Calling  & 60.0 & 46.7 & 26.7 & 11.1 &  4.4 \\
\bottomrule
\end{tabular}
\end{table}

Table~\ref{tab:intervention-resistor} reveals several key findings:

\paragraph{1. Success traces monotonically improve performance.}
Injecting even a single successful reasoning step ($n_s{=}1$) raises the ReAct success rate from 73.3\% to 86.7\% (+13.4~pp) and Tool Calling from 60.0\% to 68.9\% (+8.9~pp).
As more steps are injected, performance climbs steadily: at $n_s{=}2$, ReAct reaches 95.6\% and Tool Calling 82.2\%.
At the highest dosage ($n_s{=}n{-}1$), both agents achieve perfect 100\% success.
This monotonic improvement suggests that partial successful traces provide a strong inductive bias that the agent can effectively leverage to complete the remaining reasoning, even when the initial approach varies across sampled traces.

\paragraph{2. Failure traces monotonically degrade performance.}
Even a single failed reasoning step substantially degrades performance: ReAct drops by 24.4~pp and Tool Calling by 13.3~pp at $n_s{=}1$.
The degradation intensifies with more injected steps.
At the highest dosage ($n_s{=}n{-}1$), ReAct falls to 11.1\% and Tool Calling to just 4.4\%---a near-complete collapse.
This demonstrates that agents are unable to recover from flawed reasoning prefixes; instead, they tend to continue along the failed trajectory, amplifying rather than correcting initial errors.

\paragraph{3. The effect is asymmetric.}
At low dosages, the magnitude of failure-induced degradation exceeds the success-induced improvement.
At $n_s{=}1$, the failure penalty for ReAct ($-$24.4~pp) is nearly double the success bonus (+13.4~pp).
This asymmetry suggests that agents are more sensitive to misleading context than to helpful scaffolding---a single wrong step is harder to overcome than a single right step is to build upon.
This finding has implications for systems that retrieve prior agent traces as few-shot examples: the cost of retrieving a poor example may outweigh the benefit of retrieving a good one.

\paragraph{4. Tool Calling agents are more vulnerable to failure traces.}
While both agents converge to 100\% under full success intervention, Tool Calling degrades more severely under failure traces: 4.4\% vs.\ 11.1\% at $n_s{=}n{-}1$.
This may reflect the structured nature of tool-calling interactions, where the rigid function-call format creates stronger path dependencies---once the agent commits to a particular sequence of structured tool invocations, deviating from that established pattern becomes more difficult than in the free-form ReAct paradigm.

\paragraph{5. Implications.}
These results demonstrate that conversational context acts as a powerful prior for LLM-based agents: correct partial trajectories can nearly guarantee task success, while incorrect ones can be catastrophic.
This has practical implications for several common agent design patterns:
\begin{itemize}
    \item \textbf{Conversation history management}: Agents that maintain and carry forward conversation histories across sessions are susceptible to error propagation from failed sub-conversations.
    \item \textbf{Retrieval-augmented agent reasoning}: Systems that retrieve prior agent traces as in-context examples must carefully filter for quality, as the penalty for including a failed trace substantially exceeds the benefit of including a successful one.
    \item \textbf{Multi-agent collaboration}: When agents share reasoning artifacts, a single agent's failure trajectory can cascade through the system, degrading the performance of downstream agents that consume these artifacts.
\end{itemize}
